\documentclass[14pt,a4paper]{extarticle}

\usepackage{fontspec}
\usepackage[ukrainian,english]{babel}
\usepackage{amsmath}
\usepackage{amssymb}
\usepackage{geometry}
\usepackage{enumitem}
\usepackage{indentfirst}
\usepackage{microtype}
\usepackage{ragged2e}
\usepackage{setspace}

\geometry{
  left=25mm,
  right=10mm,
  top=20mm,
  bottom=20mm
}

\setmainfont{Times New Roman}
\setlength{\parindent}{1.25cm}
\setlength{\parskip}{0pt}
\setlength{\emergencystretch}{3em}
\onehalfspacing
\setlist{nosep,leftmargin=1.25cm,labelsep=0.25cm}

\newcommand{\eng}[1]{\foreignlanguage{english}{#1}}
\newcommand{\sectiontitle}[1]{%
  \par\vspace{\baselineskip}%
  \noindent\textbf{#1}\par
}
\newcommand{\assignmenttitle}[1]{%
  \begin{center}
  \textbf{#1}
  \end{center}
}
\newcommand{\studentline}[1]{\noindent\textbf{#1}\par}
\newcommand{\researchtopic}[1]{\noindent\textbf{Тема дисертаційного дослідження: «#1».}\par}
\newenvironment{numberedreferences}{%
  \par\vspace{\baselineskip}
  \begin{center}
  \textbf{Список використаних джерел}
  \end{center}
  \begingroup
  \setlength{\parindent}{0pt}
  \setlength{\leftskip}{0pt}
  \setlength{\rightskip}{0pt}
  \setlength{\parskip}{0pt}
  \justifying
  \begin{enumerate}[leftmargin=0.75cm,labelwidth=1cm,labelsep=0cm,align=left,label={[\arabic*]},itemsep=0pt,topsep=0pt,parsep=0pt]
}{%
  \end{enumerate}
  \endgroup
}

\begin{document}

\assignmenttitle{Назва дисципліни або роботи}

\studentline{Виконав: аспірант гр. 121-24а, Прізвище Ім'я}

\researchtopic{Методи та програмні засоби рельєфного текстурування для систем формування тривимірних графічних зображень}

\sectiontitle{1. Назва першої секції}

Текст першого абзацу. Українські університетські роботи оформлюються шрифтом Times New Roman, 14 пунктів, з міжрядковим інтервалом 1.5. Англомовний технічний термін можна залишати поруч у дужках, наприклад рельєфне текстурування (\eng{relief texturing}).

Текст другого абзацу без додаткового порожнього рядка у фінальному PDF. Вихідний LaTeX-файл може мати порожній рядок між абзацами, але вертикальний інтервал контролюється через \verb|\parskip|.

\sectiontitle{2. Формули}

Приклад математичного запису:
\[
H(u,v), \quad H(u,v) \in [0,1].
\]

Приклад моделі:
\[
D(u,v)=1-H(u,v).
\]

\sectiontitle{Висновок}

Короткий підсумок роботи.

\begin{numberedreferences}
  \item Author, A. A. Title of work. \textit{Journal or Publisher}, 2024.
  \item Author, B. B., \& Author, C. C. Title of article. \textit{Journal Title}, vol. 12, no. 3, pp. 45--67, 2023.
\end{numberedreferences}

\end{document}
